% Praca inżynierska: Opracowanie i wykonanie układu sterującego klawiaturą mechaniczną
% Kompilacja: latexmk (konfiguracja w latexmkrc), wynik w build/main.pdf
%
% UWAGA: format strony, czcionka i strona tytułowa są TYMCZASOWE.
% Po pobraniu szablonu WEFiM PWr podmienić preambułę i title.tex.

\documentclass[a4paper,12pt,oneside]{report}

% --- Język i kodowanie ---
\usepackage[T1]{fontenc}
\usepackage[utf8]{inputenc}
\usepackage[polish]{babel}
\DeclareUnicodeCharacter{03A9}{\ensuremath{\Omega}}
\DeclareUnicodeCharacter{00B5}{\ensuremath{\mu}}
\usepackage{lmodern}
\usepackage[autostyle]{csquotes}

% --- Układ strony (tymczasowo, do zgodności z szablonem) ---
\usepackage[a4paper,top=2.5cm,bottom=2.5cm,inner=3.5cm,outer=2.5cm]{geometry}
\usepackage[onehalfspacing]{setspace}
\usepackage{microtype}

% --- Matematyka, jednostki, tabele, rysunki ---
\usepackage{amsmath,amssymb}
\usepackage{siunitx}
\sisetup{output-decimal-marker={,}, group-separator={\,}, per-mode=symbol, range-phrase={\,--\,}}
\usepackage{graphicx}
\graphicspath{{figures/}}
\usepackage{booktabs,tabularx,multirow}
\usepackage{caption,subcaption}
\captionsetup{font=small,labelfont=bf}
\usepackage{enumitem}
\setlist{itemsep=0.2em}
\usepackage{xcolor}
\usepackage{listings}
\lstset{basicstyle=\ttfamily\small,breaklines=true,frame=single,columns=fullflexible}

% --- Notatki robocze (przed złożeniem: \usepackage[disable]{todonotes}) ---
\usepackage[colorinlistoftodos,textsize=footnotesize]{todonotes}
\setuptodonotes{inline}

% --- Notatki o zawartości rozdziałów (szare ramki). Ukrycie: \roboczafalse ---
\usepackage{environ}
\newif\ifrobocza \roboczatrue
\NewEnviron{zawartosc}{\ifrobocza\par\noindent\fbox{\begin{minipage}{\dimexpr\linewidth-2\fboxsep-2\fboxrule}\small\color{gray!60!black}%
  \textbf{Co opisujemy w tym rozdziale:}\begin{itemize}[leftmargin=1.5em,itemsep=0.1em,topsep=0.3em]\BODY\end{itemize}\end{minipage}}\par\medskip\fi}

% --- Bibliografia ---
\usepackage[backend=biber,style=ieee,sorting=none]{biblatex}
\addbibresource{bib/bibliografia.bib}

% --- Odnośniki (zawsze na końcu preambuły) ---
\usepackage[hidelinks,unicode]{hyperref}

% --- Skróty projektu ---
\newcommand{\nrf}{nRF52840}
\newcommand{\npm}{nPM1300}
\newcommand{\zmk}{ZMK}

\begin{document}

% TYMCZASOWA strona tytułowa: zastąpić stroną z szablonu WEFiM PWr.
\begin{titlepage}
  \centering
  {\large Politechnika Wrocławska\\Wydział Elektroniki, Fotoniki i Mikrosystemów\par}
  \vspace{1em}
  {Kierunek: Inteligentna Elektronika\par}
  \vfill
  {\Large\textsc{Praca dyplomowa inżynierska}\par}
  \vspace{2em}
  {\LARGE\bfseries Opracowanie i wykonanie układu sterującego klawiaturą mechaniczną\par}
  \vspace{1em}
  {\large\itshape Design and construction of a control system for a mechanical keyboard\par}
  \vfill
  \begin{flushright}
    Autor: [imię i nazwisko]\\[0.5em]
    Opiekun pracy: [tytuł, imię i nazwisko]
  \end{flushright}
  \vfill
  {Wrocław 2026\par}
\end{titlepage}


\pagenumbering{roman}
% Streszczenie i abstract: pisać na końcu. Ok. 200 słów każde.
% Cel, co powstało (controller MB + DB, firmware ZMK, projekt membrany), 2-3 liczby z rozdz. 7.
\chapter*{Streszczenie}
\addcontentsline{toc}{chapter}{Streszczenie}
\todo{Napisać na końcu.}

\noindent\textbf{Słowa kluczowe:} klawiatura mechaniczna, IBM Model M, Bluetooth Low Energy, nRF52840, nPM1300, ZMK

\chapter*{Abstract}
\addcontentsline{toc}{chapter}{Abstract}
\todo{Write at the end.}

\noindent\textbf{Keywords:} mechanical keyboard, IBM Model M, Bluetooth Low Energy, nRF52840, nPM1300, ZMK

\tableofcontents
\listoftodos[Do zrobienia (lista robocza, usunąć przed złożeniem)]
\cleardoublepage
\pagenumbering{arabic}

\chapter{Wstęp}\label{ch:wstep}
% Ok. 2-3 strony. Pisać pod koniec, gdy reszta jest gotowa.

\section{Motywacja, cel i zakres pracy}\label{sec:cel}
% - Model M jako klasyk (buckling spring); brak bezprzewodowości i USB-C w oryginale.
% - Założenie (notatnik 28.10): pełna kompatybilność z oryginalnym mechanizmem, drop-in replacement.
% - Cel jednym zdaniem + cele szczegółowe: projekt elektroniki, wykonanie PCB, firmware, weryfikacja pomiarowa.
% - Zakres: membrana FPC jako projekt, smart features (J3) jako kierunek rozwoju.
\todo{Napisać.}

\section{Układ pracy}\label{sec:uklad}
% - Akapit: co jest w kolejnych rozdziałach.
\todo{Napisać.}


\chapter{Analiza problemu i przegląd rozwiązań}\label{ch:analiza}
% Ok. 8-10 stron. Wymaga zebrania źródeł i zdjęć z rozbiórki.

\section{Konstrukcja IBM Model M SSK}\label{sec:modelm}
% - Buckling spring over membrane, stalowa płyta, układ SSK, oryginalny kontroler.
% - Membrana: dwie warstwy przewodzące + spacer; ogony; bolt-mod. Zdjęcia własne.
% - Źródło: patent IBM US 4,118,611 (buckling spring).
\todo{Napisać.}

\section{Matryca membranowa: ghosting, masking, NKRO}\label{sec:matryca}
% - Skanowanie wiersz/kolumna: 6 + 17 = 23 GPIO zamiast 88 wejść.
% - Rysunek prostokąta 2x2: ghosting, masking; dioda na klawisz daje NKRO.
% - Oryginalna membrana bez diod -> decyzja o własnej membranie.
\todo{Napisać.}

\section{Istniejące kontrolery zastępcze i ich porównanie}\label{sec:kontrolery}
% - Oryginalny kontroler, otwarte zamienniki (Teensy/ATmega32U4), FSSK (reference/fssk-v1.00).
% - Tabela: MCU, przewód/bezprzewodowo, NKRO, ingerencja w obudowę, otwartość.
% - Wniosek: brak bezprzewodowego drop-in z NKRO.
\todo{Napisać.}

\section{Łączność bezprzewodowa w klawiaturach}\label{sec:ble}
% - BLE HID over GATT, interwał połączenia vs opóźnienie vs prąd; USB HID; dongle 2,4 GHz krótko.
\todo{Napisać.}

\section{Oprogramowanie klawiatur: QMK, ZMK}\label{sec:qmkzmk}
% - QMK vs ZMK (Zephyr, BLE, devicetree, MIT); uzasadnienie wyboru: HD 'Firmware: ZMK na Zephyr 4.1'.
\todo{Napisać.}


\chapter{Wymagania i koncepcja systemu}\label{ch:wymagania}
% Ok. 5-7 stron. Materiał: HD 'Założenia architektoniczne'.
\begin{zawartosc}
  \item Lista wymagań funkcjonalnych i niefunkcjonalnych w tabeli z numerami (weryfikowane w rozdz. 7).
  \item Architektura systemu na schemacie blokowym: płytka główna (MB), daughterboard z USB-C (DB), dwie folie membrany, ogniwo Li-Po.
  \item Wybór kluczowych układów (nRF52840, nPM1300) z porównaniem alternatyw i uzasadnieniem: bare chip zamiast modułu, PMIC zamiast osobnej ładowarki i regulatora.
  \item Wstępne bilanse: liczba GPIO, budżet prądowy szyny 3,3 V, szacowany czas pracy z danych katalogowych.
\end{zawartosc}

\section{Wymagania funkcjonalne oraz niefunkcjonalne}\label{sec:wymagania}
% - Tabela W1..Wn, każde wymaganie weryfikowane w sekcji 7.7.
% - Funkcjonalne: BLE + USB HID, NKRO, ładowanie z USB-C, dioda statusu, SWD, raport baterii.
% - Niefunkcjonalne: drop-in, czas pracy (wartość do ustalenia), produkcja JLC, koszt, rozszerzalność (J3).
\todo{Napisać.}

\section{Architektura systemu}\label{sec:architektura}
% - Schemat blokowy: MB, DB, łącze 1x10, membrana (J4, J7), ogniwo (J2). Główny rysunek pracy.
% - Podział na płytki i projekty KiCada (CL A2).
\todo{Napisać.}

\section{Wybór układów}\label{sec:wybor}
% - MCU: nRF52840 bare chip vs moduł; tabela alternatyw.
% - PMIC: nPM1300 vs ładowarka + LDO/buck; tabela alternatyw.
% - Ogniwo i złącze (CL A26, A27, P11).
% - Budżet: GPIO (HD 'Przydział GPIO'), prąd szyny 3,3 V (HD 'Obciążenie szyny 3,3 V'), szacunek czasu pracy.
\todo{Napisać.}


\chapter{Projekt układu elektronicznego}\label{ch:elektronika}
% Ok. 18-22 stron. Rytm każdej sekcji: wymagania bloku -> wycinek schematu -> dobór elementów
% -> obliczenia ze wzorem i referencją do datasheetu -> co weryfikuje rozdz. 7.
\begin{zawartosc}
  \item Zasilanie: nPM1300 (BUCK2 3,3 V, ładowarka CC/CV, LOADSW1, tryb Ship), ochrona przed odwrotną polaryzacją; ogniwo bez NTC (R13 0 Ω zamiast termistora) i konsekwencje tej decyzji.
  \item Mikrokontroler nRF52840: konfiguracja zasilania, odsprzęganie, kwarce z obliczeniem pojemności obciążenia, programowanie SWD, przydział GPIO.
  \item Tor radiowy: dopasowanie anteny, linia 50 Ω (CPWG) na wybranym stackupie, antena w obudowie ze stalową płytą.
  \item Interfejsy: USB-C z zabezpieczeniami, łącze MB-DB, I2C, złącze rozszerzeń J3, dioda statusu.
  \item Membrana FPC: matryca 6 x 17 col2row, folia ROW (2 warstwy, z diodami) i folia COL (1 warstwa) na oryginalnym spacerze, złącza J4/J7.
  \item Projekt PCB (4 warstwy, rozmieszczenie, pętle prądowe, punkty pomiarowe) oraz BOM, koszt i produkcja w JLC (controller zamówiony 16.11, membrana 23.11).
\end{zawartosc}

\section{PMIC nPM1300}\label{sec:pmic}
% - BUCK2 = 3,3 V (VSET2 470k), BUCK1 off, VDDIO, I2C, GPIO0 przerwanie, GPIO1 LOADSW1.
% - Ładowarka CC/CV, bez NTC (R13 0 Ω, firmware thermistor-ohms = 0), tryb Ship, ochrona przed odwrotną polaryzacją.
% - Obliczenia: HD 'Blok PMIC: obliczenia' (tętnienia cewki, zakres regulacji, straty ładowarki, spadek na Q1).
% - Tabela połączeń: HD 'Blok PMIC: połączenia'.
\todo{Napisać.}

\section{MCU nRF52840}\label{sec:mcu}
% - Konfiguracja 5 (VDD = VDDH, REG1 DC/DC), odsprzęganie, HFXO/LFXO, SWD TC2030, RESET, SWO.
% - Obliczenia: HD 'Obliczenia: obciążenie kwarców' (historia błędnego CL, CL M1).
% - Przydział GPIO: HD 'Przydział GPIO', CL M7.
\todo{Napisać.}

\section{Tor RF}\label{sec:rf}
% - Dopasowanie 4-elementowe z config 7, U.FL, antena w obudowie ze stalową płytą.
% - CPWG 50 Ohm na JLC04121H-3313: 0,135/0,15 mm (CL A30, HD 'PCB: stackup'); L1 Murata (CL A19).
\todo{Napisać.}

\section{Interfejsy}\label{sec:interfejsy}
% - USB-C bez 5,1k, ESD, koordynacja zabezpieczeń VBUS (HD).
% - Łącze MB-DB 1x10 i pinout; I2C pull-upy (HD 'Obliczenie: pull-upy I2C').
% - J3 (HD 'Smart features'), dioda WS2812B przez LOADSW1 (HD 'Dioda statusu').
\todo{Napisać.}

\section{Matryca klawiszy i membrana FPC}\label{sec:membrana}
% - Matryca 6x17, col2row, dioda na klawisz, J4 (18) i J7 (6) (HD 'Złącza matrycy').
% - Folia ROW 2L z diodami + folia COL 1L na oryginalnym spacerze (hardware/membrane, zamówione 23.11, production/membrane).
\todo{Napisać.}

\section{Projekt PCB}\label{sec:pcb}
% - Stackup 4L 1,2 mm (CL A24), rozmieszczenie, pętla BUCK2 i NT1 (CL A22), via-in-pad POFV (A18, A25).
% - Reguły DRC (A29), TP i R3 (A15, A28), DB wyłamywany. Widoki warstw, render 3D.
\todo{Napisać.}

\section{BOM, koszt i produkcja}\label{sec:bom}
% - JLC Basic vs Extended (HD 'BOM pod montaż JLC'), części lutowane ręcznie (U2, J3, J5).
% - Zamówienie rev A (CL A31), tabela kosztów.
\todo{Napisać.}


\chapter{Konstrukcja mechaniczna i montaż}\label{ch:mechanika}
% Wersja skrócona, ok. 2-4 strony. Jeśli materiału za mało: przenieść do 4.8 'Integracja mechaniczna'.

\section{Umieszczenie w obudowie}\label{sec:obudowa}
% - Model Onshape, położenie MB i DB, gniazdo USB. Zrzuty z Onshape.
\todo{Napisać.}

\section{Uchwyt daughterboardu}\label{sec:uchwytdb}
% - Koncepcja; wykonanie po obronie.
\todo{Napisać.}

\section{Koszyk baterii}\label{sec:koszyk}
% - Koncepcja v5 (bateria/koszyk-wymiary.md), otwory w PCB (CL A27).
\todo{Napisać.}

\section{Montaż membrany i bolt-mod}\label{sec:boltmod}
% - Procedura i stan.
\todo{Napisać.}


\chapter{Oprogramowanie}\label{ch:firmware}
% Ok. 5-7 stron. Repo zmk-config (CL F5): płytka ibm_m_ssk, boards/blonski3/ibm_m_ssk.
\begin{zawartosc}
  \item Firmware ZMK na Zephyrze 4.1: własna definicja płytki (HWMv2), devicetree, skanowanie matrycy col2row 17 x 6, tryby USB i BLE.
  \item Obsługa nPM1300: parametry ładowania, regulator, LOADSW1, przerwanie; praca bez NTC.
  \item Raportowanie stanu baterii i sterowanie diodą statusu.
  \item Wgrywanie firmware'u (SWD, ewentualnie bootloader UF2) i sposób aktualizacji.
\end{zawartosc}

\section{ZMK na Zephyr 4.1}\label{sec:zmk}
% - Definicja płytki HWMv2, devicetree, kscan col2row (piny CL M7), BLE + USB, NKRO.
\todo{Napisać.}

\section{Obsługa nPM1300}\label{sec:fwpmic}
% - Ładowarka 400 mA / 4,20 V, limit VBUS 500 mA, BUCK2, NTC, LOADSW1, przerwanie (HD 'Ustawienia startowe').
\todo{Napisać.}

\section{Obsługa własnych modułów}\label{sec:moduly}
% - Raport baterii, dioda statusu; zakres zależny od czasu.
\todo{Napisać.}

\section{Bootloader UF2 i aktualizacje}\label{sec:uf2}
% - Do 18.12 wgrywanie przez SWD; UF2 jako stan/kierunek rozwoju.
\todo{Napisać.}


\chapter{Uruchomienie i badania prototypu}\label{ch:badania}
% Ok. 8-12 stron. Plan pomiarów: praca/plan-pracy.md sekcja 3 (9 pozycji).

\section{Stanowisko pomiarowe i metodyka}\label{sec:metodyka}
% - Sprzęt, warunki (J-Link odłączony, J3 pusty, DIN niski), niepewności.
\todo{Napisać.}

\section{Uruchamianie krok po kroku}\label{sec:uruchomienie}
% - Lutowanie U2, Z1 (BUCK2 bez nRF, R3 wylutowany), R1 SWD, pierwszy firmware.
\todo{Napisać.}

\section{Zasilanie i ładowanie}\label{sec:pomzasilanie}
% - Z1, Z2, Z4b, Ł1, Ł3, Ł5.
\todo{Napisać.}

\section{Radio}\label{sec:pomradio}
% - R1, R5 (RSSI w funkcji odległości, z płytą i bez).
\todo{Napisać.}

\section{Funkcje klawiatury}\label{sec:pomklawiatura}
% - K1-K2 na makiecie matrycy.
\todo{Napisać.}

\section{Pobór energii i czas pracy}\label{sec:pomenergia}
% - E1, E2.
\todo{Napisać.}

\section{Weryfikacja wymagań}\label{sec:weryfikacja}
% - Tabela W1..Wn -> test -> wynik -> spełnione/niespełnione.
\todo{Napisać.}


\chapter{Podsumowanie i wnioski}\label{ch:podsumowanie}
% Ok. 2-3 strony. Pisać na końcu.

\section{Realizacja celu i wymagań}\label{sec:realizacja}
% - Odwołanie do tabeli w sekcji 7.7.
\todo{Napisać.}

\section{Napotkane problemy}\label{sec:problemy}
% - Błędny CL kwarcu (M1), odwrócony tranzystor ochronny, BUCK1 vs VSET, WS2812C niedostępna,
% - R10 47k zamiast 470k w BOM, nPM1300 bez stanu w JLC, zmiany złącza ogniwa (A21, A26).
\todo{Napisać.}

\section{Rozwój projektu}\label{sec:rozwoj}
% - Membrana i montaż w obudowie, UF2, moduły J3, dynamiczny limit VBUS z CC, rev B.
\todo{Napisać.}



\nocite{*} % tymczasowo: pokazuje całą bazę; usunąć, gdy pojawią się cytowania
\printbibliography[heading=bibintoc,title={Bibliografia}]
\listoffigures
\listoftables

\appendix
% Załączniki. Schemat i widoki PCB jako PDF z KiCada (wstawiać przez \includegraphics lub pdfpages).
\chapter{Schemat układu sterującego}\label{app:schemat}
\todo{Wstawić PDF schematu controllera.}

\chapter{Widoki płytki PCB i stackup}\label{app:pcb}
\todo{Wstawić widoki warstw.}

\chapter{Wykaz elementów (BOM)}\label{app:bom}
\todo{Tabela BOM.}

\chapter{Konfiguracja ZMK}\label{app:zmk}
\todo{Listingi devicetree.}

\chapter{Punkty pomiarowe}\label{app:tp}
\todo{Tabela TP i R3 (stan po A28).}


\end{document}
